% GENERATED FILE. Edit canonical lessons, then run npm run generate:books.
% canonical-source-hash: fnv1a64:7ac3ff63ed55c7dd
% canonical-lessons: HI-C137-svetar, HI-C137-anda, HI-C137-gost, HI-C137-alu, HI-C137-pyaz

\chapter{Jumper, Egg, Meat, Potato, Onion}
\label{ch:hi-jumper-egg-meat-potato-onion}
\hlchaptermodality{137}

\begin{chapteropening}
\textbf{By the end of this chapter:} \emph{I can say a jumper, an egg, meat, a potato and an onion.}

Complete the last lesson of chapter 137: I can say a jumper, an egg, meat, a potato and an onion.
\end{chapteropening}

\section[\texorpdfstring{\hi{स्वेटर} (sveṭar)}{स्वेटर (sveṭar)}]{\hi{स्वेटर} (sveṭar) — a jumper}

\label{lesson:HI-C137-svetar}



\begin{quote}
Before the new one: say the Hindi for a cap, then the Hindi for a coat.
\end{quote}

\subsection*{You'll want to know: \hi{स्वेटर}}
\textbf{\hi{स्वेटर}} — \emph{sveṭar} — \textquotedblleft{}a jumper\textquotedblright{}.

\begin{grammarlens}[title={What you've built}]
One more word for this chapter.
\end{grammarlens}

\subsection*{Your turn}
\begin{itemize}
\raggedright
  \item \emph{Say it:} \emph{sveṭar}
  \item \emph{Say it:} \emph{sveṭar}, once more
  \item \emph{Say it:} say \emph{sveṭar}
  \item \emph{Recall:} say the Hindi for a cap, then the Hindi for a coat, then say \emph{sveṭar} again
\end{itemize}

\begin{tcolorbox}[breakable,colback=teal!4,colframe=teal!35!black,title={Before you move on}]
What does \hi{स्वेटर} mean? (\textquotedblleft{}A jumper\textquotedblright{}.) Say it once more.
\end{tcolorbox}
\section[\texorpdfstring{\hi{अंडा} (aṇḍā)}{अंडा (aṇḍā)}]{\hi{अंडा} (aṇḍā) — an egg}

\label{lesson:HI-C137-anda}



\begin{quote}
Before the new one: say the Hindi for a coat, then the Hindi for a jumper.

Greet the day through: \textbf{\hi{सुप्रभात}}, \textbf{\hi{शुभ} \hi{दोपहर}}, \textbf{\hi{शुभ} \hi{संध्या}}. Then say \textbf{\hi{सुबह} \hi{से} \hi{शाम} \hi{तक}}, from \textbf{\hi{सुबह}} to \textbf{\hi{शाम}}.
\end{quote}

\subsection*{You'll want to know: \hi{अंडा}}
\textbf{\hi{अंडा}} — \emph{aṇḍā} — \textquotedblleft{}an egg\textquotedblright{}.

\begin{grammarlens}[title={What you've built}]
One more word for this chapter.
\end{grammarlens}

\subsection*{Your turn}
\begin{itemize}
\raggedright
  \item \emph{Say it:} \emph{aṇḍā}
  \item \emph{Say it:} \emph{aṇḍā}, once more
  \item \emph{Say it:} say \emph{aṇḍā}
  \item \emph{Recall:} say the Hindi for a coat, then the Hindi for a jumper, then say \emph{aṇḍā} again
\end{itemize}

\begin{tcolorbox}[breakable,colback=teal!4,colframe=teal!35!black,title={Before you move on}]
What does \hi{अंडा} mean? (\textquotedblleft{}An egg\textquotedblright{}.) Say it once more.
\end{tcolorbox}
\section[\texorpdfstring{\hi{गोश्त} (gośt)}{गोश्त (gośt)}]{\hi{गोश्त} (gośt) — meat}

\label{lesson:HI-C137-gost}



\begin{quote}
Before the new one: say the Hindi for a jumper, then the Hindi for an egg.
\end{quote}

\subsection*{You'll want to know: \hi{गोश्त}}
\textbf{\hi{गोश्त}} — \emph{gośt} — \textquotedblleft{}meat\textquotedblright{}.

\begin{grammarlens}[title={What you've built}]
One more word for this chapter.
\end{grammarlens}

\subsection*{Your turn}
\begin{itemize}
\raggedright
  \item \emph{Say it:} \emph{gośt}
  \item \emph{Say it:} \emph{gośt}, once more
  \item \emph{Say it:} say \emph{gośt}
  \item \emph{Recall:} say the Hindi for a jumper, then the Hindi for an egg, then say \emph{gośt} again
\end{itemize}

\begin{tcolorbox}[breakable,colback=teal!4,colframe=teal!35!black,title={Before you move on}]
What does \hi{गोश्त} mean? (\textquotedblleft{}Meat\textquotedblright{}.) Say it once more.
\end{tcolorbox}
\section[\texorpdfstring{\hi{आलू} (ālū)}{आलू (ālū)}]{\hi{आलू} (ālū) — a potato}

\label{lesson:HI-C137-alu}



\begin{quote}
Before the new one: say the Hindi for an egg, then the Hindi for meat.
\end{quote}

\subsection*{You'll want to know: \hi{आलू}}
\textbf{\hi{आलू}} — \emph{ālū} — \textquotedblleft{}a potato\textquotedblright{}.

\begin{grammarlens}[title={What you've built}]
One more word for this chapter.
\end{grammarlens}

\subsection*{Your turn}
\begin{itemize}
\raggedright
  \item \emph{Say it:} \emph{ālū}
  \item \emph{Say it:} \emph{ālū}, once more
  \item \emph{Say it:} say \emph{ālū}
  \item \emph{Recall:} say the Hindi for an egg, then the Hindi for meat, then say \emph{ālū} again
\end{itemize}

\begin{tcolorbox}[breakable,colback=teal!4,colframe=teal!35!black,title={Before you move on}]
What does \hi{आलू} mean? (\textquotedblleft{}A potato\textquotedblright{}.) Say it once more.
\end{tcolorbox}
\section[\texorpdfstring{\hi{प्याज़} (pyāz)}{प्याज़ (pyāz)}]{\hi{प्याज़} (pyāz) — an onion}

\label{lesson:HI-C137-pyaz}



\begin{quote}
Before the new one: say the Hindi for meat, then the Hindi for a potato.
\end{quote}

\subsection*{You'll want to know: \hi{प्याज़}}
\textbf{\hi{प्याज़}} — \emph{pyāz} — \textquotedblleft{}an onion\textquotedblright{}.

\begin{grammarlens}[title={What you've built}]
One more word for this chapter.
\end{grammarlens}

\subsection*{Your turn}
\begin{itemize}
\raggedright
  \item \emph{Say it:} \emph{pyāz}
  \item \emph{Say it:} \emph{pyāz}, once more
  \item \emph{Say it:} say \emph{pyāz}
  \item \emph{Recall:} say the Hindi for meat, then the Hindi for a potato, then say \emph{pyāz} again
\end{itemize}

\begin{tcolorbox}[breakable,colback=teal!4,colframe=teal!35!black,title={Before you move on}]
What does \hi{प्याज़} mean? (\textquotedblleft{}An onion\textquotedblright{}.) Say it once more.
\end{tcolorbox}
